\documentclass[11pt]{article}
\usepackage{fontspec}
\setmainfont[
  Path=fonts/,
  Extension=.ttf,
  UprightFont=*-400,
  BoldFont=*-700
]{NotoSerifSC}
\setsansfont[
  Path=fonts/,
  Extension=.ttf,
  UprightFont=*-400,
  BoldFont=*-700
]{NotoSerifSC}
\XeTeXlinebreaklocale "zh"
\XeTeXlinebreakskip=0pt plus 1pt
\usepackage[a4paper,margin=1.70cm]{geometry}
\usepackage{amsmath,amssymb,mathtools,bm}
\usepackage{booktabs,longtable,array}
\usepackage{enumitem}
\setlength{\parindent}{2em}
\setlength{\parskip}{0.35em}
\setlist{nosep,leftmargin=2em}
\allowdisplaybreaks[3]
\numberwithin{equation}{section}
\pagestyle{plain}

\newcommand{\R}{\mathbb{R}}
\newcommand{\E}{\mathbb{E}}
\newcommand{\Pp}{\mathbb{P}}
\newcommand{\1}{\mathbf{1}}
\newcommand{\T}{\mathsf{T}}
\newcommand{\F}{\mathrm{F}}
\newcommand{\cN}{\mathcal{N}}
\newcommand{\cL}{\mathcal{L}}
\newcommand{\cS}{\mathcal{S}}
\newcommand{\cH}{\mathcal{H}}
\newcommand{\cA}{\mathcal{A}}
\newcommand{\cM}{\mathcal{M}}
\newcommand{\cT}{\mathcal{T}}
\newcommand{\diag}{\operatorname{diag}}
\newcommand{\tr}{\operatorname{tr}}
\newcommand{\Var}{\operatorname{Var}}
\newcommand{\Cov}{\operatorname{Cov}}
\newcommand{\softmax}{\operatorname{softmax}}
\newcommand{\KL}{D_{\mathrm{KL}}}
\newcommand{\sg}{\operatorname{sg}}
\begin{document}
    \begin{center}
      {\LARGE\bfseries 44\_Mapping Networks}
    \end{center}
    \vspace{0.8em}

    \section{符号与维度}
    \begin{longtable}{>{$}l<{$} >{$}l<{$} p{10cm}}
    \toprule
    \text{符号} & \text{维度} & \text{含义}\\
    \midrule
    \theta & \R^P & 目标网络完整参数。\\
        z & \R^d & 低维潜参数，$d\ll P$。\\
        g_\psi & \R^d\to\R^P & 映射网络。\\
        \cM & -- & 假设中的低维参数流形。\\
        L_\theta,L_\ell & \R_{>0} & 模型和损失的 Lipschitz 常数。\\
    \bottomrule
    \end{longtable}

    所有向量默认是列向量；未特别说明时，期望同时对数据、噪声和采样时刻取值。下文只展开论文中承担方法逻辑的公式，并把所需假设写在推导发生的位置。

\section{局部坐标图与映射定理的严格构造}

设 $\mathcal M\subset\R^P$ 是 $d^*$ 维 $C^2$ 嵌入流形，
$\theta^*\in\mathcal M$。按定义，存在开集 $U\subset\R^{d^*}$、
$0\in U$ 和局部坐标图
\begin{equation}
 \phi:U\to V\subset\mathcal M,\qquad
 \phi(0)=\theta^*,\qquad \operatorname{rank}J_\phi(u)=d^*.
\end{equation}
若任务损失在 $B(\theta^*,r)$ 内满足
\begin{equation}
 |L(\theta_1)-L(\theta_2)|
 \le L_\ell L_\theta\|\theta_1-\theta_2\|_2,
\end{equation}
给定 $0<\epsilon\le L_\ell L_\theta r$，令
\begin{equation}
 \delta=\frac{\epsilon}{L_\ell L_\theta}\le r.
\end{equation}
由 $\phi$ 在 $0$ 连续，存在 $\eta>0$ 使
\begin{equation}
 \|u\|<\eta\Rightarrow
 \|\phi(u)-\theta^*\|<\delta.
\end{equation}
于是任取该邻域内 $z^*$，
\begin{equation}
 |L(\phi(z^*))-L(\theta^*)|
 \le L_\ell L_\theta\|\phi(z^*)-\theta^*\|<\epsilon.
\end{equation}

论文希望得到定义在整个 $\R^{d^*}$ 上的 $C^2$ 映射。严格做法是先选
$U_1\Subset U_2\Subset U$，$0\in U_1$，再取光滑 cutoff
$\psi\in C_c^\infty(U_2)$ 且 $\psi=1$ 于 $U_1$。定义
\begin{equation}
 g(u)=
 \begin{cases}
 \theta^*+\psi(u)[\phi(u)-\theta^*],&u\in U_2,\\
 \theta^*,&u\notin U_2.
 \end{cases}
\end{equation}
因为 $\psi$ 在 $U_2$ 边界附近恒为零，两段及其任意阶导数平滑拼接，
$g\in C^2(\R^{d^*},\R^P)$。这避免直接在 $U$ 外书写未定义的 $\phi(u)$。

注意取 $z^*=0$ 就有
\begin{equation}
 g(0)=\phi(0)=\theta^*,
\end{equation}
损失误差严格为零。因此这个存在性定理本身来自“已知最优点位于流形并以它为坐标图中心”，
并没有证明某个预先随机固定的网络能找到或表示该坐标图。

\section{复合 Lipschitz 常数的来源}

若对所有输入 $x$，
\begin{equation}
 \|f_{\theta_1}(x)-f_{\theta_2}(x)\|
 \le L_\theta\|\theta_1-\theta_2\|,
\end{equation}
且逐样本损失对预测满足
\begin{equation}
 |\ell(y_1,y)-\ell(y_2,y)|\le L_\ell\|y_1-y_2\|,
\end{equation}
则期望任务损失
$L(\theta)=\E_{x,y}\ell(f_\theta(x),y)$ 满足
\begin{align}
 |L(\theta_1)-L(\theta_2)|
 &\le\E|\ell(f_{\theta_1}(x),y)-\ell(f_{\theta_2}(x),y)|\\
 &\le L_\ell\E\|f_{\theta_1}(x)-f_{\theta_2}(x)\|\\
 &\le L_\ell L_\theta\|\theta_1-\theta_2\|.
\end{align}
若交叉熵的预测概率可趋近零，它在概率坐标中没有全局有限 Lipschitz 常数；
需要把 logits 或概率限制在紧集，故论文的结论是局部而非全局。

\section{潜变量梯度是切空间拉回}

令 $\theta=g_\psi(z)$，任务梯度 $q=\nabla_\theta L\in\R^P$。
链式法则
\begin{equation}
 \boxed{\nabla_zL(g(z))=J_g(z)^Tq.}
\end{equation}
潜空间一步 $z^+=z-\eta J_g^Tq$ 在参数空间的一阶变化
\begin{equation}
 \Delta\theta
 \approx J_g(z)(z^+-z)
 =-\eta J_gJ_g^Tq.
\end{equation}
若 $J_g=U\Sigma V^T$，则
\begin{equation}
 J_gJ_g^T=U\Sigma^2U^T.
\end{equation}
所以更新只在切空间 $\operatorname{col}(J_g)$ 内，并按奇异值平方预条件。
它通常不是欧氏正交投影；正交投影为
\begin{equation}
 P_T=J_g(J_g^TJ_g)^{-1}J_g^T
\end{equation}
（满列秩时）。若希望参数空间沿 $-P_Tq$，潜更新应使用自然梯度
\begin{equation}
 \Delta z=-\eta(J_g^TJ_g)^{-1}J_g^Tq.
\end{equation}

奇异值界给出
\begin{equation}
 \sigma_{\min}\|P_Tq\|
 \le\|J_g^Tq\|
 \le\sigma_{\max}\|P_Tq\|.
\end{equation}
小 $\sigma_{\min}$ 造成某些切向方向学习极慢，条件数
$\kappa(J_g)=\sigma_{\max}/\sigma_{\min}$ 是低维优化是否良态的关键。

\section{局部可解性的逆函数推导}

设初点 $z_0$、$\theta_0=g(z_0)$，残差
$r=\theta^*-\theta_0$。Taylor 展开
\begin{equation}
 g(z_0+\Delta z)=\theta_0+J\Delta z+R(\Delta z),
 \qquad\|R(\Delta z)\|\le\frac M2\|\Delta z\|^2,
\end{equation}
其中 $J=J_g(z_0)$，$M$ 控制二阶导数。
若 $r$ 位于 $\operatorname{col}(J)$，且 $J$ 满列秩，可取
\begin{equation}
 \Delta z=J^\dagger r,\qquad
 \|\Delta z\|\le\frac{\|r\|}{\sigma_{\min}(J)}.
\end{equation}
线性项精确抵消残差，剩余
\begin{equation}
 \|g(z_0+\Delta z)-\theta^*\|
 \le\frac M{2\sigma_{\min}(J)^2}\|r\|^2.
\end{equation}
这给出论文 $O(\|r\|^2)$ 声称所需的明确条件。

若 $r$ 含法向分量 $r_\perp=(I-P_T)r\ne0$，任何一阶潜更新都无法消除它：
\begin{equation}
 \min_{\Delta z}\|J\Delta z-r\|=|r_\perp\|.
\end{equation}
因此仅假设 $\|r\|$ 小并不足以保证二次剩余；还需目标位于局部映射像，或法向残差本身为二阶小量。

\section{加性权重调制的导数}

映射网络权重由潜变量调制
\begin{equation}
 \omega(z)=\omega_0+Bz,\qquad
 g(z)=g_{\omega(z)}(z).
\end{equation}
因为 $z$ 同时作为网络输入和权重来源，全导数是
\begin{equation}
 \boxed{\frac{dg}{dz}
 =\frac{\partial g_\omega(z)}{\partial z}
 +\frac{\partial g_\omega(z)}{\partial\omega}B.}
\end{equation}
只计算第一项会漏掉调制路径。对逐元素规则
$\omega_{ij}\leftarrow\omega_{ij}+\alpha z_i$，矩阵 $B$ 是稀疏复制算子；
同一 $z_i$ 调制多个权重，因此其梯度累加这些连接的贡献。

加性调制仍是一个特定参数化。要声称它满足映射定理，必须证明其像包含定理构造的局部坐标图，
或至少以所需精度逼近；仅有 $C^2$ 光滑性并不保证像经过任意 $\theta^*$。

\section{稳定、平滑和对齐正则的精确含义}

稳定损失
\begin{equation}
 L_{\rm stab}=\E_{\epsilon\sim\cN(0,\sigma^2I)}
 \|F(z+\epsilon)-F(z)\|_2^2
\end{equation}
在小 $\sigma$ 下 Taylor 展开
\begin{equation}
 F(z+\epsilon)-F(z)=J_F(z)\epsilon+O(\|\epsilon\|^2).
\end{equation}
故
\begin{align}
 L_{\rm stab}
 &=\E\epsilon^TJ_F^TJ_F\epsilon+O(\sigma^3)\\
 &=\tr(J_F^TJ_F\E\epsilon\epsilon^T)+O(\sigma^3)\\
 &=\sigma^2\|J_F(z)\|_F^2+O(\sigma^3).
\end{align}
它在小噪声极限等价于 Jacobian 正则。

论文平滑项
\begin{equation}
 L_{\rm smooth}=\|J_{M_\phi}(z)\|_F^2
\end{equation}
只惩罚一阶导数幅值，并不单独“保证 $C^2$”。若网络激活本身不是 $C^2$（例如 ReLU），
函数仍可能只分段光滑；要控制曲率还需 Hessian 项
\begin{equation}
 \|\nabla_z^2M_\phi(z)\|_F^2.
\end{equation}

余弦对齐
\begin{equation}
 L_{\rm align}=1-\frac{z^Tw}{\|z\|\|w\|}
\end{equation}
对 $z$ 的梯度
\begin{equation}
 \nabla_zL_{\rm align}
 =-\frac1{\|z\|\|w\|}
 \left(w-\frac{z^Tw}{\|z\|^2}z\right),
\end{equation}
它只改变方向且在 $z=0$ 奇异，实际实现需范数下界。

\section{全局拓扑限制}

单一坐标图只能覆盖与 $\R^d$ 开集微分同胚的局部区域。若最优参数集合是圆环、球面或多个不连通分支，
不存在一张无奇点、单射的欧氏全局坐标图覆盖全部集合。
连续 $g:\R^d\to\R^P$ 可以覆盖或穿过这些区域，但可能非单射，且不同潜点映射到同一参数。
因此流形假设支持局部降维，不自动推出单个潜向量参数化整个训练轨迹的全局可学习性。

\section{光滑损失的下降引理}

若 $L$ 的梯度是 $\beta$-Lipschitz：
\begin{equation}
 \|\nabla L(x)-\nabla L(y)\|
 \le\beta\|x-y\|,
\end{equation}
则沿线段积分
\begin{align}
 L(y)-L(x)
 &=\int_0^1\nabla L(x+t(y-x))^T(y-x)dt\\
 &=\nabla L(x)^T(y-x)+R,
\end{align}
余项满足
\begin{align}
 |R|
 &\le\int_0^1\beta t\|y-x\|^2dt
 =\frac\beta2\|y-x\|^2.
\end{align}
故下降引理
\begin{equation}
 L(y)\le L(x)+\nabla L(x)^T(y-x)
 +\frac\beta2\|y-x\|^2.
\end{equation}
取梯度步 $y=x-\eta\nabla L(x)$：
\begin{equation}
 L(y)\le L(x)-\eta\left(1-\frac{\beta\eta}{2}\right)
 \|\nabla L(x)\|^2.
\end{equation}
当 $0<\eta<2/\beta$ 时每个精确梯度步下降；深网中 $\beta$ 未知且随位置变化，
这仍解释了学习率过大造成不稳定的局部机制。

\section{二阶近似与曲率方向}

在 $x$ 附近 Taylor 展开
\begin{equation}
 L(x+\Delta)=L(x)+g^T\Delta
 +\frac12\Delta^TH\Delta+O(\|\Delta\|^3).
\end{equation}
若 $H=U\diag(\lambda_i)U^T$，在本征方向 $u_i$ 上步长 $a$ 的变化
\begin{equation}
 \Delta L\approx a g_i+\frac12\lambda_i a^2.
\end{equation}
$\lambda_i>0$ 为局部碗形，$\lambda_i<0$ 是负曲率逃离鞍点方向。
加阻尼的 Newton 步
\begin{equation}
 \Delta=-(H+\lambda I)^{-1}g
\end{equation}
在特征方向缩放为 $-g_i/(\lambda_i+\lambda)$，避免接近零或负特征值造成巨大步长。

\section{链式 Jacobian 与条件数}

复合映射 $y=f_L\circ\cdots\circ f_1(x)$ 的 Jacobian
\begin{equation}
 J_{y\leftarrow x}=J_LJ_{L-1}\cdots J_1.
\end{equation}
谱范数上界
\begin{equation}
 \|J_{y\leftarrow x}\|_2\le\prod_{\ell=1}^{L}\|J_\ell\|_2.
\end{equation}
若各层最小奇异值存在，
\begin{equation}
 \sigma_{\min}(J_{y\leftarrow x})
 \ge\prod_\ell\sigma_{\min}(J_\ell).
\end{equation}
因此条件数至多乘法累积：
\begin{equation}
 \kappa(J_{y\leftarrow x})
 \le\prod_\ell\kappa(J_\ell).
\end{equation}
低维映射、递归模块和物理传播都面临同一问题：表示存在不代表反向梯度数值良态。

\section{随机梯度的偏差与方差}

设样本梯度 $g(\theta;\xi)$ 满足
\begin{equation}
 \E_\xi g(\theta;\xi)=\nabla L(\theta),
 \qquad \E\|g-\nabla L\|^2\le\sigma^2.
\end{equation}
独立批量 $B$ 的平均
\begin{equation}
 \bar g_B=\frac1B\sum_{i=1}^{B}g_i
\end{equation}
满足
\begin{equation}
 \E\bar g_B=\nabla L,\qquad
 \E\|\bar g_B-\nabla L\|^2\le\frac{\sigma^2}{B}.
\end{equation}
若样本相关，协方差项出现：
\begin{equation}
 \Var(\bar g_B)=\frac1{B^2}\left[
 \sum_i\Var(g_i)+\sum_{i\ne j}\Cov(g_i,g_j)
 \right].
\end{equation}
扩大批量只有在交叉协方差不主导时才近似按 $1/B$ 降噪。

若使用有偏估计 $\E\widehat g=\nabla L+b$，均方误差分解
\begin{equation}
 \E\|\widehat g-\nabla L\|^2
 =\|b\|^2+\tr\Cov(\widehat g).
\end{equation}
直通估计、停止梯度目标和截断反传通常以可控偏差换低方差或低内存。

\section{Adam 的尺度归一化}

逐坐标 Adam 更新
\begin{align}
 m_t&=\beta_1m_{t-1}+(1-\beta_1)g_t,\\
 v_t&=\beta_2v_{t-1}+(1-\beta_2)g_t^2,\\
 \widehat m_t&=m_t/(1-\beta_1^t),\\
 \widehat v_t&=v_t/(1-\beta_2^t),\\
 \theta_{t+1}&=\theta_t-\eta
 \frac{\widehat m_t}{\sqrt{\widehat v_t}+\epsilon}.
\end{align}
若梯度长期乘常数 $c>0$，忽略 $\epsilon$ 时
$\widehat m$ 乘 $c$、$\sqrt{\widehat v}$ 也乘 $c$，更新近似不变。
但瞬态、符号变化和 $\epsilon$ 会破坏精确尺度不变性。

AdamW 把权重衰减解耦：
\begin{equation}
 \theta_{t+1}=(1-\eta\lambda)\theta_t
 -\eta\frac{\widehat m_t}{\sqrt{\widehat v_t}+\epsilon}.
\end{equation}
这不同于把 $\lambda\theta$ 加入梯度后再经过自适应分母；后者会按坐标二阶矩缩放正则化。

\section{范数约束与投影梯度}

约束集合 $\mathcal C$ 上的投影梯度
\begin{equation}
 \theta^+=\Pi_{\mathcal C}(\theta-\eta g),
 \qquad
 \Pi_{\mathcal C}(x)=\arg\min_{y\in\mathcal C}\|x-y\|^2.
\end{equation}
对欧氏球 $\mathcal C=\{x:\|x\|\le R\}$，
\begin{equation}
 \Pi_{\mathcal C}(x)=
 \begin{cases}
 x,&\|x\|\le R,\\
 Rx/\|x\|,&\|x\|>R.
 \end{cases}
\end{equation}
对流形约束，局部一阶步先投影到切空间，再用 retraction 回到流形；
仅把参数写成 $g(z)$ 相当于用参数化隐式满足约束，但会引入 $J_g^TJ_g$ 的坐标度量。

\section{Lipschitz 误差的层间传播}

真实模块 $f_\ell$、近似模块 $\widehat f_\ell$ 满足
\begin{equation}
 \|f_\ell(x)-\widehat f_\ell(x)\|\le\epsilon_\ell,
 \qquad
 \|f_\ell(x)-f_\ell(y)\|\le L_\ell\|x-y\|.
\end{equation}
令真实、近似中间状态差 $e_\ell$。三角不等式给出
\begin{align}
 e_{\ell+1}
 &=\|f_\ell(h_\ell)-\widehat f_\ell(\widehat h_\ell)\|\\
 &\le L_\ell e_\ell+\epsilon_\ell.
\end{align}
递归展开
\begin{equation}
 e_L\le
 \left(\prod_{j=0}^{L-1}L_j\right)e_0
 +\sum_{\ell=0}^{L-1}
 \left(\prod_{j=\ell+1}^{L-1}L_j\right)\epsilon_\ell.
\end{equation}
早期层误差乘更多后续 Lipschitz 因子，所以量化、低秩近似或局部训练的相同单层误差，
放在不同深度并不等价。

\section{复杂度必须区分四个量}

参数量、FLOPs、激活内存和通信量是不同指标。以线性层
$Y=XW$、$X\in\R^{B\times n}$、$W\in\R^{n\times m}$ 为例：
\begin{align}
 \text{参数量}&=nm,\\
 \text{前向乘加量}&\asymp Bnm,\\
 \text{主要激活量}&=B(n+m),\\
 \text{数据并行梯度通信}&\asymp nm.
\end{align}
低秩因子可同时减少参数与乘法，但动态稀疏若实现仍做稠密 kernel，参数/FLOPs 公式不会自动转化为墙钟加速。
冻结参数减少梯度和优化器状态，却不一定减少前向计算或输入梯度计算。

\end{document}
